\section*{ID6538: Numerical Recipe (Practical: How to Simulate \& Visualize)}
\paragraph{Metadata.}
PiID: 4028308711232 | RTEID: RTE:P25388.5059:T4.1333 | UUID: 9913bbd9-2e94-5b8a-a3ae-cad1ac57e49f

\paragraph{Canonical Statement.}
**Key parameter suggestions (Earth-like scale):** - \textbackslash{}(R=6.371\textbackslash{}times10\textasciicircum{}6\textbackslash{} \textbackslash{}mathrm\{m\}\textbackslash{}), - \textbackslash{}(\textbackslash{}Omega=7.2921\textbackslash{}times10\textasciicircum{}\{-5\}\textbackslash{} \textbackslash{}mathrm\{rad/s\}\textbackslash{}), - shallow-water equivalent depth \textbackslash{}(H\textbackslash{}sim 100\textbackslash{})–\textbackslash{}(1000\textbackslash{}) m if you want some gravity-wave coupling, - small viscosity \textbackslash{}(\textbackslash{}nu\textbackslash{}sim 10\textasciicircum{}4\textbackslash{})–\textbackslash{}(10\textasciicircum{}6\textbackslash{} \textbackslash{}mathrm\{m\textasciicircum{}2/s\}\textbackslash{}) or use hyperdiffusion.

\paragraph{Canonical Equation.}
\[
R=6.371\times10^6\ \mathrm{m}
\]

\paragraph{Dictionary.}
Numerical Recipe (Practical: How to Simulate \& Visualize) — R=6.371×10\textasciicircum{}6 m

\paragraph{Encyclopedia.}
Numerical Recipe (Practical: How to Simulate \& Visualize) is an algebraic definition, written R=6.371×10\textasciicircum{}6 m. It is an algebraic definition expressing the quantity directly in terms of the variables on its right-hand side. Within the Trawin Zero Point registry it serves as one element of the framework's mathematical narrative.
